\chapter{Anomalous And Topological Transport}
\label{ch:thermal-anomalous-topological-transport}

\ConventionsMixed

This chapter studies transport terms fixed by anomalies and topology.  The
logic is real-time response first: a transport coefficient is defined by a
correlator or by the response of a generating functional to background
sources, and anomaly equations constrain the allowed terms.

\section{Hydrodynamic Source Datum}

Let \(W[g,A]\) be the equilibrium generating functional for a QFT with
stress tensor \(T^{\mu\nu}\) and current \(J^\mu\).  Variations define
\[
  \delta W
  =
  \frac12\int\dd^Dx\,\sqrt{-g}\,
  \langle T^{\mu\nu}\rangle\,\delta g_{\mu\nu}
  +
  \int\dd^Dx\,\sqrt{-g}\,
  \langle J^\mu\rangle\,\delta A_\mu .
\]
Hydrodynamic variables are a temperature \(T\), velocity \(u^\mu\), and
chemical potential \(\mu\), with
\[
  u^\mu u_\mu=-1.
\]
The ideal constitutive relations are
\[
  T^{\mu\nu}_{(0)}=\epsilon u^\mu u^\nu+p\Delta^{\mu\nu},\qquad
  J^\mu_{(0)}=n u^\mu,
\]
where
\[
  \Delta^{\mu\nu}=g^{\mu\nu}+u^\mu u^\nu .
\]

\section{Anomaly Equation}

In four dimensions, a \(U(1)^3\) anomaly has the local equation
\[
  \nabla_\mu J^\mu
  =
  \frac{C}{32\pi^2}
  \epsilon^{\mu\nu\rho\sigma}F_{\mu\nu}F_{\rho\sigma}.
\]
The coefficient \(C\) is part of the current normalization.  The magnetic
field and vorticity in the local rest frame are
\[
  B^\mu=\frac12\epsilon^{\mu\nu\rho\sigma}u_\nu F_{\rho\sigma},
  \qquad
  \omega^\mu=\frac12\epsilon^{\mu\nu\rho\sigma}
  u_\nu\nabla_\rho u_\sigma .
\]
The first-derivative parity-odd current terms have the form
\[
  J^\mu_{\mathrm{odd}}
  =
  \xi_B B^\mu+\xi_\omega\omega^\mu .
\]
The coefficients \(\xi_B,\xi_\omega\) are transport coefficients.  They must
be defined either by Kubo formulae or by differentiating an equilibrium
generating functional.

For the explicit four-dimensional formulas below, let \(J_R\) and \(J_L\) be
the source-normalized right- and left-handed number currents of one massless
Dirac fermion, and set
\[
  J_V=J_R+J_L,\qquad J_A=J_R-J_L,\qquad
  \mu_R=\mu_V+\mu_A,\qquad \mu_L=\mu_V-\mu_A .
\]
The vector source is \(A_V\).  In the convention of
Chapter~\ref{ch:volume-ii-20}, preserving the vector Ward identity gives
\[
  \nabla_\mu J_A^\mu
  =
  \frac1{16\pi^2}
  \epsilon^{\mu\nu\rho\sigma}
  (F_V)_{\mu\nu}(F_V)_{\rho\sigma}
  \qquad
  \hbox{for a unit source charge.}
\]
If a physical electromagnetic source \(A_{\rm em}\) couples with charge
\(Q=eq\), then \(A_V=Q A_{\rm em}\).  The current conjugate to
\(A_{\rm em}\) is \(QJ_V\), and the magnetic response acquires the factor
\(Q^2=e^2q^2\).

\section{Kubo Definitions}

In flat space at temperature \(T\), define retarded correlators
\[
  G^{R}_{AB}(t,\mathbf x)
  =
  -\ii\theta(t)\langle[A(t,\mathbf x),B(0)]\rangle_\beta .
\]
For a current response to a weak magnetic field, the chiral magnetic
coefficient is
\[
  \xi_B
  =
  \lim_{k\to0}
  \frac{\ii}{2k}
  \epsilon_{ijn}\,
  G^R_{J^iJ^j}(\omega=0,\mathbf k=k\hat e_n).
\]
This formula is a definition of the transport coefficient in the declared
static, homogeneous limit.  If the limits \(\omega\to0\) and \(k\to0\) fail
to commute, the chosen order is part of the observable.

\section{Equilibrium Derivation}

Place the system on a stationary background
\[
  \dd s^2= -\ee^{2\sigma}(dt+a_i dx^i)^2+h_{ij}dx^i dx^j,
  \qquad
  A=A_0(dt+a_i dx^i)+A_i dx^i .
\]
Gauge and Kaluza--Klein invariant three-dimensional fields are
\[
  \mathcal A_i=A_i-A_0a_i,\qquad A_0 .
\]
An anomalous equilibrium functional contains Chern--Simons-type terms
\[
  W_{\mathrm{odd}}
  =
  \int_{\Sigma_3}
  \left[
    c_1 A_0\,\mathcal A\wedge\dd\mathcal A
    +
    c_2 A_0^2\,\mathcal A\wedge\dd a
    +
    c_3 A_0^3\,a\wedge\dd a
  \right],
\]
with coefficients constrained by the anomaly and by large gauge
transformations.  Differentiating this functional with respect to
\(\mathcal A_i\) and \(a_i\) gives the parity-odd equilibrium pieces of
\(J^i\) and \(T^{0i}\).  The derivation is algebraic once the equilibrium
functional and its gauge-variation convention are fixed.

\section{Chiral Magnetic Coefficient}

The anomaly fixes a Chern--Simons term in the three-dimensional equilibrium
functional.  Hold the axial source constant in Euclidean time,
\[
  (A_A)_0=\mu_A,\qquad \partial_i\mu_A=0,
\]
and set the spatial axial field to zero.  The inflow functional for the
vector-vector-axial anomaly contains the five-dimensional local term
\[
  W_5[A_A,A_V]
  =
  -\frac1{4\pi^2}
  \int_{X_5} A_A\wedge F_V\wedge F_V ,
\]
whose axial gauge variation is
\[
  \delta_\beta W_5
  =
  -\frac1{4\pi^2}\int_{\partial X_5}\beta\,F_V\wedge F_V .
\]
Since
\[
  F_V\wedge F_V
  =
  \frac14\epsilon^{\mu\nu\rho\sigma}
  (F_V)_{\mu\nu}(F_V)_{\rho\sigma}\,\dd^4x ,
\]
this is the inflow form of the axial anomaly displayed above, with the sign
fixed by the convention that the combined bulk-boundary system is invariant.
Reducing the same local term on the thermal circle gives the parity-odd
equilibrium functional
\[
  W_{\rm CME}[A_V;\mu_A]
  =
  \frac{\mu_A}{4\pi^2}
  \int_{\Sigma_3} A_V\wedge\dd A_V .
\]
In components,
\[
  \int_{\Sigma_3} A_V\wedge\dd A_V
  =
  \int\dd^3x\,
  \epsilon^{ijk}(A_V)_i\partial_j(A_V)_k .
\]
Varying with respect to the spatial vector source gives
\[
  \frac{\delta W_{\rm CME}}{\delta (A_V)_i}
  =
  \frac{\mu_A}{2\pi^2}
  \epsilon^{ijk}\partial_j(A_V)_k
  =
  \frac{\mu_A}{2\pi^2}B_V^i .
\]
Thus the covariant equilibrium chiral magnetic response of the
source-normalized vector current is
\begin{equation}
  J_V^i
  =
  \frac{\mu_A}{2\pi^2}B_V^i .
  \label{eq:cme-source-normalized}
\end{equation}
For the physical electromagnetic current of a Dirac fermion of charge
\(Q=eq\), \(B_V=Q B_{\rm em}\) and \(J_{\rm em}=QJ_V\), hence
\begin{equation}
  J_{\rm em}^i
  =
  \frac{e^2q^2\mu_A}{2\pi^2}B_{\rm em}^i .
  \label{eq:cme-electromagnetic-current}
\end{equation}

The same coefficient is obtained without invoking hydrodynamics by counting
lowest Landau levels.  For \(B_V=B\hat z\), the degeneracy per unit transverse
area is \(|B|/(2\pi)\).  The right-handed lowest Landau level has group
velocity along \(B\), while the left-handed lowest Landau level has group
velocity opposite to \(B\).  Filling momenta from \(0\) to \(\mu_R\) and
\(0\) to \(\mu_L\) gives
\[
  J_R^z=\frac{\mu_R}{4\pi^2}B,\qquad
  J_L^z=-\frac{\mu_L}{4\pi^2}B .
\]
Therefore
\[
  J_V^z=J_R^z+J_L^z
  =
  \frac{\mu_R-\mu_L}{4\pi^2}B
  =
  \frac{\mu_A}{2\pi^2}B,
\]
which agrees with \eqref{eq:cme-source-normalized}.  This derivation also
shows why the order of limits in the Kubo formula is part of the observable:
the current is an equilibrium response to a static magnetic source and a
specified axial chemical potential.

\section{Chiral Vortical Coefficients}

For one massless Dirac fermion, the covariant first-derivative vortical
currents in the same source-normalized convention are
\begin{align}
  J_V^\mu\big|_\omega
  &=
  \frac{\mu_V\mu_A}{\pi^2}\,\omega^\mu ,
  \label{eq:vector-cve-dirac}\\
  J_A^\mu\big|_\omega
  &=
  \left[
    \frac{\mu_V^2+\mu_A^2}{2\pi^2}
    +\frac{T^2}{6}
  \right]\omega^\mu .
  \label{eq:axial-cve-dirac}
\end{align}
These formulas follow by adding the two Weyl contributions
\[
  J_R^\mu\big|_\omega
  =
  \left(\frac{\mu_R^2}{4\pi^2}+\frac{T^2}{12}\right)\omega^\mu,
  \qquad
  J_L^\mu\big|_\omega
  =
  -\left(\frac{\mu_L^2}{4\pi^2}+\frac{T^2}{12}\right)\omega^\mu .
\]
The vector current is the sum \(J_R+J_L\), while the axial current is the
difference \(J_R-J_L\).  Hence the \(T^2/12\) terms cancel in the vector
current and add in the axial current.  There is therefore no universal
source-normalized vector-current term proportional to \(\mu_A T\) in this
Dirac example; the universal temperature-squared contribution belongs to the
axial vortical current and is tied to the mixed axial-gravitational anomaly.

\section{Entropy Constraint}

Let
\[
  \nabla_\mu T^{\mu\nu}=F^{\nu\lambda}J_\lambda,\qquad
  \nabla_\mu J^\mu=\mathcal A
\]
with anomaly density \(\mathcal A\).  The canonical entropy current is
\[
  S^\mu_{\mathrm{can}}
  =
  p\beta^\mu-\beta_\nu T^{\mu\nu}-\alpha J^\mu,
  \qquad
  \beta^\mu=\frac{u^\mu}{T},\quad \alpha=\frac{\mu}{T}.
\]
Using the thermodynamic identity
\[
  \dd p=s\,\dd T+n\,\dd\mu
\]
and the conservation equations, the divergence of \(S^\mu_{\mathrm{can}}\)
contains the dissipative products plus the anomaly term
\[
  -\alpha\,\mathcal A .
\]
Parity-odd additions to \(S^\mu\) can cancel nondissipative anomaly-induced
terms, leaving nonnegative entropy production from dissipative tensors.  This
calculation fixes relations among \(\xi_B,\xi_\omega\) and anomaly
coefficients, up to terms not seen by the entropy analysis.

\section{Topological Terms And Contact Terms}

Equilibrium Chern--Simons terms may change local current correlators by
contact terms.  A transport coefficient must therefore state whether the
current is covariant or consistent, which background functional is used, and
which contact terms are included in the retarded correlator.  The physically
measured response is attached to the source protocol, not merely to a local
operator symbol.

\begin{openproblem}[Nonperturbative anomalous transport]
\label{op:nonperturbative-anomalous-transport}
Derive anomaly-induced transport coefficients in interacting QFT directly
from nonperturbative KMS states, local current algebras, Ward identities,
contact terms, and hydrodynamic scaling limits, with the order of limits and
source protocol stated as part of the theorem.
\end{openproblem}
